% Source: https://artofproblemsolving.com/community/c3h2483605_nt_problems
% Stable identity: aops-topic-2483605-q10
The square roots require $-\dfrac34\le x\le\dfrac12$. Put $u=\sqrt{1-2x}\ge0$ and $v=\sqrt{4x+3}\ge0$. Then
\[2u^2+v^2=2-4x+4x+3=5\qquad\text{and}\qquad 2u+v=\sqrt{15}.\]
Substituting $v=\sqrt{15}-2u$ into the first equation:
\[2u^2+15-4\sqrt{15}\,u+4u^2=5\iff6u^2-4\sqrt{15}\,u+10=0\iff\left(\sqrt6\,u-\sqrt{10}\right)^2=0.\]
So $u=\sqrt{\dfrac{10}{6}}=\sqrt{\dfrac53}$, and then $v=\sqrt{15}-2\sqrt{\dfrac53}=\sqrt{\dfrac53}\,(3-2)=\sqrt{\dfrac53}\ge0$. From $1-2x=u^2=\dfrac53$ we get $x=-\dfrac13$, which also satisfies $4x+3=\dfrac53=v^2$.

Check: $2\sqrt{\dfrac53}+\sqrt{\dfrac53}=3\sqrt{\dfrac53}=\sqrt{15}$. The unique solution is $\boxed{x=-\dfrac13}$.

(Equivalently, by the Cauchy--Schwarz inequality $(2u+v)^2\le(2+1)\left(2u^2+v^2\right)=15$, with equality only when $u=v$.)
